\subsubsection{Descomposición armónica del retorno electrónico}
\label{sec:el-espectro-retorno}

La trayectoria APP--TRIT--TPK produce primero el registro
\(\tau_e\). La representación armónica actúa después sobre ese
registro y explicita la información que evalúa \(K_\Omega\).
En esta realización, \(\pi\) y la unidad compleja son las
coordenadas ya construidas; los caracteres cíclicos constituyen un
lenguaje de lectura de la palabra producida.

Para \(\tau\in\mathbb R^{12}\), con índices en
\(\mathbb Z/12\mathbb Z\), fijamos
\begin{equation}
 T_k(\tau)=\sum_{j=0}^{11}\tau_j
             \exp\!\left(-\frac{2\pi i kj}{12}\right),\qquad
 \theta_k=\frac{2\pi k}{12}.
 \label{eq:el-transformada-finita}
\end{equation}
La fórmula extiende linealmente la representación del registro
trítico al espacio real de palabras. Las correlaciones \(C_s\),
las ventanas \(A_\ell\) y el lector \(P_6\) conservan las
definiciones anteriores en ese dominio ampliado.

\begin{theorem}[Representación espectral del lector de retorno]
\label{thm:el-lector-espectral}
Para toda palabra real de longitud doce,
\begin{align}
 C_s(\tau)&=\frac1{12}\sum_{k=0}^{11}|T_k(\tau)|^2
                         \cos(s\theta_k),
 \label{eq:el-correlacion-espectral}\\
 \Omega(\tau)&=\frac1{12}\sum_{k=0}^{11}|T_k(\tau)|^2w_k,
 \label{eq:el-omega-espectral}\\
 w_k&=-2\cos\theta_k-4\cos(2\theta_k)-6\cos(3\theta_k),\notag\\
 P_6(\tau)&=\frac1{24}\sum_{k=0}^{11}|T_k(\tau)|^2(-1)^k.
 \label{eq:el-p6-espectral}
\end{align}
\end{theorem}
\begin{proof}
Al desarrollar el módulo cuadrado se obtiene
\[
 |T_k|^2=\sum_{a,b=0}^{11}\tau_a\tau_b
           \exp\!\left(-\frac{2\pi i k(a-b)}{12}\right).
\]
La identidad de ortogonalidad
\[
 \frac1{12}\sum_{k=0}^{11}
       \exp\!\left(\frac{2\pi i k n}{12}\right)
 =\begin{cases}1,&n\equiv0\pmod{12},\\0,&n\not\equiv0\pmod{12}
   \end{cases}
\]
se deduce sumando una progresión geométrica. Al multiplicar la
expresión de \(|T_k|^2\) por \(\exp(is\theta_k)\) y sumar en
\(k\), sólo permanecen los pares con \(a-b\equiv s\).
La suma resultante es \(C_s\); al tomar partes reales resulta
\eqref{eq:el-correlacion-espectral}. La combinación
\(-2C_1-4C_2-6C_3\), ya deducida de las ventanas, proporciona
\eqref{eq:el-omega-espectral}. Finalmente,
\(C_6=2P_6\) y \(\cos(6\theta_k)=(-1)^k\) dan
\eqref{eq:el-p6-espectral}.
\end{proof}

\begin{corollary}[Tres modos del registro central]
\label{cor:el-tres-modos}
La palabra \(\tau_e=(1,1,1,-1,-1,-1)^2\) tiene las potencias
espectrales
\begin{equation}
 |T_2|^2=64,\qquad |T_6|^2=16,\qquad |T_{10}|^2=64,
 \label{eq:el-potencias-espectrales}
\end{equation}
y las restantes potencias son cero. Su evaluación produce
\begin{equation}
 \begin{aligned}
 \Omega(\tau_e)&=\frac{64\cdot7+16\cdot4+64\cdot7}{12}=80,\\
 P_6(\tau_e)&=\frac{64+16+64}{24}=6.
 \end{aligned}
 \label{eq:el-evaluacion-espectral}
\end{equation}
\end{corollary}
\begin{proof}
Con \(z_k=\exp(-i\theta_k)\), la repetición de la sextupla da
\[
 T_k=(1+(-1)^k)(1-z_k^3)(1+z_k+z_k^2).
\]
El primer factor anula los índices impares; el segundo anula
\(k=0,4,8\). En los tres índices restantes se obtiene
\[
 T_2=4-4i\sqrt3,\qquad T_6=4,\qquad T_{10}=4+4i\sqrt3.
\]
Sus módulos cuadrados son los indicados. La definición de los
pesos da \(w_2=w_{10}=7\) y \(w_6=4\); todos esos índices son
pares. Las fórmulas del teorema concluyen la evaluación.
\end{proof}

La sustitución en el carácter de retorno proporciona una segunda
expresión exacta de su dependencia respecto de la historia:
\begin{equation}
 \kappa_e=
 \frac1{12}\sum_{k=0}^{11}|T_k(\tau_e)|^2
       \left(w_k+\frac{\Delta_4}{2}(-1)^k\right)-54\alpha.
 \label{eq:el-kappa-espectral}
\end{equation}
El término \(-54\alpha\) conserva la procedencia cronológica
de \(K_\Omega\). Los otros dos términos evalúan las potencias
armónicas con los pesos de las ventanas y la correlación
antipodal. La composición electrónica completa
\eqref{eq:el-formula-completa} recibe así una lectura explícita
de sus modos de retorno. La norma \(\sqrt3/4\), la transferencia
\(\varphi/\pi^2\), los censos \(22=27-5\) y \(15=5\cdot3\),
el defecto \(\Delta_4\) y el cociente \(R_{\mathrm{act}}\)
conservan las construcciones previamente demostradas y reunidas en
la tabla~\ref{tab:el-factores}.

\paragraph{Orden, origen de fase e información conservada.}
La palabra alternante \(\tau_{\rm alt}=(1,-1)^6\) contiene los
mismos seis signos positivos y seis negativos que \(\tau_e\).
Su transformada tiene únicamente \(T_6=12\), de modo que
\(\Omega(\tau_{\rm alt})=144\cdot4/12=48\).
Este testigo prueba la sensibilidad al orden de un lector que
el histograma de signos deja indeterminado. Su función es delimitar
el dominio del lector; la selección de una especie requiere la
sección y la realización correspondientes.

Si \((S_a\tau)_j=\tau_{j+a}\), entonces
\(T_k(S_a\tau)=\exp(ia\theta_k)T_k(\tau)\).
Las potencias y el carácter de retorno conservan, por ello, las
traslaciones cíclicas. Los coeficientes complejos completos
reconstruyen la palabra mediante
\[
 \tau_j=\frac1{12}\sum_{k=0}^{11}T_k(\tau)\exp(ij\theta_k),
\]
mientras sus módulos cuadrados conservan la autocorrelación.
La coordenatización integral distingue, además, las historias de eventos de
\eqref{eq:el-historias-centrales}. Esta jerarquía precisa qué
información transmite cada lectura del mismo transporte.

Las identidades anteriores poseen demostraciones algebraicas para
todo el dominio declarado. Su comprobación reproducible evalúa los
testigos en \(\mathbb Z[z]/(z^4-z^2+1)\), comprueba las doce
traslaciones y recorre las \(4096\) palabras binarias de longitud
doce para contrastar la identidad entre ventanas y correlaciones.
Los tres modos son componentes del registro dodecafásico; las
realizaciones de estados espaciales ligados utilizan, adicionalmente,
los operadores dinámicos y de interacción de su propio dominio.
